% Estilo LaTeX para el tutorial SBVI
\usepackage{booktabs}
\usepackage{longtable}
\usepackage{array}
\usepackage{tabularx}
\usepackage{ragged2e}
\usepackage{microtype}
\usepackage{fancyhdr}
\usepackage[most]{tcolorbox}
\usepackage{enumitem}
\usepackage{tikz}
\usetikzlibrary{arrows.meta,positioning,shapes.geometric}
\usepackage{siunitx}
\usepackage{xcolor}
\usepackage{titlesec}
\usepackage{hyperref}

\definecolor{BioBlue}{HTML}{17365D}
\definecolor{BioTeal}{HTML}{1F6E78}
\definecolor{BioLight}{HTML}{EEF4F7}
\definecolor{BioGray}{HTML}{5C6770}
\definecolor{BioWarn}{HTML}{8A4B08}
\definecolor{BioRed}{HTML}{8B1E2D}
\definecolor{BioGreen}{HTML}{2F6B4F}

\hypersetup{
  colorlinks=true,
  linkcolor=BioBlue,
  urlcolor=BioTeal,
  citecolor=BioGreen
}

\pagestyle{fancy}
\fancyhf{}
\fancyhead[L]{\small Señales Bioeléctricas y Visualización Interactiva}
\fancyhead[R]{\small sEMG y monitoreo post-ictus}
\fancyfoot[C]{\thepage}
\renewcommand{\headrulewidth}{0.4pt}

\titleformat{\chapter}[display]
  {\bfseries\Huge\color{BioBlue}}
  {\filleft\Large\color{BioTeal}\chaptertitlename\ \thechapter}
  {1ex}
  {\titlerule\vspace{1ex}\filright}
  [\vspace{1ex}\titlerule]

\titleformat{\section}
  {\Large\bfseries\color{BioBlue}}
  {\thesection}{0.7em}{}
\titleformat{\subsection}
  {\large\bfseries\color{BioTeal}}
  {\thesubsection}{0.7em}{}

\newtcolorbox{clinicalbox}[1][]{
  enhanced,
  breakable,
  colback=BioLight,
  colframe=BioBlue,
  boxrule=0.8pt,
  arc=2mm,
  left=2mm,right=2mm,top=1.5mm,bottom=1.5mm,
  title=Solicitud clínica,
  fonttitle=\bfseries,
  #1
}

\newtcolorbox{engineeringbox}[1][]{
  enhanced,
  breakable,
  colback=white,
  colframe=BioTeal,
  boxrule=0.8pt,
  arc=2mm,
  left=2mm,right=2mm,top=1.5mm,bottom=1.5mm,
  title=Decisión de ingeniería,
  fonttitle=\bfseries,
  #1
}

\newtcolorbox{warningbox}[1][]{
  enhanced,
  breakable,
  colback=yellow!4,
  colframe=BioWarn,
  boxrule=0.8pt,
  arc=2mm,
  left=2mm,right=2mm,top=1.5mm,bottom=1.5mm,
  title=Advertencia metodológica,
  fonttitle=\bfseries,
  #1
}

\newtcolorbox{checkpointbox}[1][]{
  enhanced,
  breakable,
  colback=green!3,
  colframe=BioGreen,
  boxrule=0.8pt,
  arc=2mm,
  left=2mm,right=2mm,top=1.5mm,bottom=1.5mm,
  title=Punto de verificación,
  fonttitle=\bfseries,
  #1
}

\setlist[itemize]{topsep=3pt,itemsep=2pt}
\setlist[enumerate]{topsep=3pt,itemsep=2pt}
\renewcommand{\arraystretch}{1.15}
